\section{Основные термины и понятния}

\paragraph{Агент.} ~-- это система, которая взаимодействует со средой, воспринимает её состояние и выбирает действия с целью достижения определённой цели или максимизации функции полезности. Формально, на каждом шаге времени $t$ агент наблюдает состояние среды $S_t$ и выбирает действие $A_t$ в соответствии со своей \emph{политикой} ~-- вероятностным распределением действий:
\begin{equation}
\pi(A_t \mid S_t) = P(A_t = a \mid S_t = s),
\end{equation}
где $\pi$ ~-- стратегия (policy) агента.  

Эволюция взаимодействия агента со средой описывается последовательностью:
\begin{equation}
S_1, A_1, R_1, S_2, A_2, R_2, \dots, S_m,
\end{equation}
где $R_t$ ~-- получаемое агентом вознаграждение на шаге $t$.  

Цель агента ~-- максимизация \emph{ожидаемого суммарного вознаграждения} за эпизод:
\begin{equation}
J(\pi) = \mathbb{E}_{\pi, T}\left[\sum_{t=1}^{m} \gamma^{t-1} R_t\right],
\end{equation}
где $\gamma \in [0,1]$ ~-- коэффициент дисконтирования, определяющий вклад будущих наград.

Для оценки качества состояния и действий агент может использовать \emph{функции ценности}:
\begin{align}
V^{\pi}(s) &= \mathbb{E}_{\pi, T}\!\left[\sum_{t=1}^{m} \gamma^{t-1} R_t \mid S_1 = s\right], \\
Q^{\pi}(s,a) &= \mathbb{E}_{\pi, T}\!\left[\sum_{t=1}^{m} \gamma^{t-1} R_t \mid S_1 = s, A_1 = a\right].
\end{align}
Оптимальная политика $\pi^*$ выбирается как:
\begin{equation}
\pi^*(a \mid s) = \arg\max_{a} Q^{\pi^*}(s,a).
\end{equation}

\paragraph{Марковский процесс принятия решений (Markov Decision Process).} Описывается пятёркой
\begin{equation}
\mathcal{M} = \langle \mathcal{S}, \mathcal{A}, T, R, \gamma \rangle,
\end{equation}
где
\begin{equation}
\begin{aligned}
& \mathcal{S}\text{~-- множество возможных состояний среды,}\\
&\mathcal{A} \text{~-- множество допустимых действий,}\\
&T(S_{t+1} \mid S_t, A_t) \text{~-- функция переходов (динамика среды),}\\
&R(S_t, A_t) \text{~-- функция вознаграждения,}\\
&\gamma \text{~-- коэффициент дисконтирования.}
\end{aligned}
\end{equation}

Динамика среды определяется распределением:
\begin{equation}
P(S_{t+1} = s' \mid S_t = s, A_t = a) = T(s' \mid s, a).
\end{equation}

Цель агента в MDP ~-- найти стратегию $\pi$, максимизирующую ожидаемую дисконтированную сумму вознаграждений:
\begin{equation}
\pi^* = \arg\max_{\pi} \mathbb{E}_{T,\pi}\!\left[\sum_{t=1}^{m} \gamma^{t-1} R(S_t, A_t)\right].
\end{equation}

Модель MDP основывается на предположениях:
\begin{equation}
\begin{aligned}
&\text{1. время дискретно;} \\
&\text{2. следующее состояние зависит только от текущего состояния и действия (свойство Маркова);} \\
&\text{3. функция переходов стационарна во времени.}
\end{aligned}
\end{equation}

Эти условия задают формальную основу для построения и анализа систем обучения с подкреплением.

\paragraph{Обучение с подкреплением (Reinforcement Learning).} Парадигма машинного обучения, в которой агент взаимодействует с окружающей средой, выполняя действия и получая награды, стремясь максимизировать суммарное вознаграждение со временем.

\paragraph{Функция вознаграждения (Reward Function).} Правило или алгоритм, присваивающее числовую награду агенту в зависимости от состояния среды и действия. В статье обозначается как $\Theta^R$ ~-- параметр определяющий поведение функции вознаграждения.

\paragraph{Вмешательство в награду (Reward Tampering).} Ситуация, когда агент учится вмешиваться в механизм получения награды, чтобы получать более высокие награды не соответствующим целям образом.

\paragraph{Диаграмма причинного влияния (Causal Influence Diagram).} Графическое представление зависимости между переменными, действиями, наблюдениями и функциями вознаграждения, помогающее выявить потенциальные уязвимости к вмешательству в награду.
